\chapter{FUNDAMENTOS DA ALEATORIEDADE EM COMPUTAÇÃO} % Capitulo 2


    \section{Entropia e Fontes de Incerteza}

        Captação de eventos físicos imprevisíveis como base para sistemas criptográficos seguros.

        Entropia é uma medida da aleatoriedade ou imprevisibilidade dos dados. 
        No contexto da criptografia, a entropia é usada para gerar 
        números aleatórios ou chaves que são essenciais para a comunicação segura e a criptografia. 
        Sem uma boa fonte de entropia, os protocolos criptográficos podem se tornar vulneráveis a ataques 
        que exploram a previsibilidade das chaves geradas.

        Na maioria dos sistemas operacionais, 
        a entropia é gerada pela coleta de eventos aleatórios de diversas fontes, 
        como interrupções de hardware, movimentos do mouse, 
        pressionamentos de teclas e atividade do disco. 
        Esses eventos são inseridos em um conjunto de entropia, 
        que é então usado para gerar números aleatórios quando necessário.

        Protocolos criptográficos, como SSL/TLS, SSH e IPsec, 
        dependem fortemente da entropia para gerar números aleatórios e chaves para criptografia, 
        autenticação e proteção de integridade.
        A aleatoriedade desses números é crucial para a segurança da comunicação. 
        Se um atacante conseguir prever os números aleatórios usados em um protocolo criptográfico, 
        ele poderá comprometer a confidencialidade, a integridade ou a autenticidade da comunicação.
        Se a chave mestra prévia for previsível ou tendenciosa, 
        um atacante poderá interceptar a mensagem criptografada, 
        descriptografá-la e obter acesso às chaves de sessão. 
        Isso permitiria ao atacante espionar a comunicação, 
        modificar os dados ou se passar pelo servidor.\nocite{6.2Entropia}

        Um tipo de entropia. É a entropia de ``Shannon''. Que utiliza umm tratamento probabilístico
        de uma função logarítmica.\cite{6.2shannonEntropia} Sendo sua formula:

        \begin{equation}
            H(X) = - \sum_{i=1}^{n} P(x_i) \log_b P(x_i)
            \label{eq:shannon_entropy}
        \end{equation}

        Em que:
        \begin{itemize}
            \item $H(X)$ é a entropia em bits;
            \item $p_i$ é a probabilidade de ocorrência do resultado $i$;
            \item $n$ é o número total de resultados possíveis.
        \end{itemize}

        Uma limitação determinística é em cima de um processador de maquina de estados finitos.
        Se acaso for utilizado as mesmas variáveis, o resultado permanece inalterado.\@ Onde se um invasor 
        tiver a certeza da seed de origem da criptografia. Ela por si mesma, não é aleatória. Para evitar tais acontecimentos
        se faz necessário que o software faça a captação da incerteza do mundo físico.

        A entropia deve ter força superior ou igual a segurança do sistema criptográfico. 
        Entropia adicional poder ser necessária para compensar a perda de entropia durante o processo de geração de chaves, mas não necessariamente adiona mais segurança.
        Ter mais entropia do que a quantidade avaliada é aceitavel, ter menos do que o recomendado pode ser fatal para a segurança do sistema criptográfico. 
        O comprimento minimo depende do mecanismo criptográfico específico e do nível de segurança desejado.
        A entrada de entropia e a seed resultante devem ser consistentes com a segurança dos dados protegidos pelo consumidor. 
        Em um modelo DRBG a seed e as chaves geradas deevem possuir o mesmo nivel de segurança, pois no modelo DRBG depende da entropia usada, seguindo o critério de segurança critico (CSP).%\nocite{NIST 800-90A, pg 18/19 8.6.3}



    \section{A Semente (Seed) como Elemento Crítico}

        Falar sobre esse ponto, pois vai ser utilizado na explicação dos protocolos tanto do funcionamento de verificação do NIST 800--22, do NIST IR 8547 e do NIST IR 7977 pos quantico 

        Uma sequência de bits usada como entrada para um gerador de números pseudoaleatórios (DRBG). A semente é a chave inicial, usada internamente para a concepção do DRBG. Sua entropia deve ser suficiente para garantir a segurança do DRBG.
        A semente é um elemento crítico, pois determina a sequência de números pseudoaleatórios gerados pelo DRBG.\@
        Se a semente for previsível ou tiver baixa entropia, o DRBG pode ser vulnerável a ataques que exploram essa fraqueza. Portanto, é essencial que a semente seja gerada a partir de uma fonte de entropia confiável.


        Uma sequência de bits usada como entrada para um mecanismo DRBG.\@ A semente determinará uma parte do interior estado do DRBG, e sua entropia deve ser suficiente para apoiar a força de segurança do DRBG
        Para obter uma seed suficientemente aleatória. Deve se submeter a uma estimativa de entropia para as saidas da fonte de ruido, baseada na propria fonte da entropia. Essa estimativa é denominada de $\mathit{H_{\text{submitter}}}$.
        Depois de possuir uma trilha de estimativas, a estimativa de entropia minima por amostra chamada de $\mathit{H_{\text{original}}}$ é determinada. 
        Se o tamanho for superior a 256 bits, ele precisara ser reduzido para menor ou igual a 256 bits.

        Se o conjunto de dados não for binário, é necessario uma entropia adicional em bit, denominada de $\mathit{H_{\text{bitstring}}}$ é estimada. O conjunto de dados sequencias possui L amostras, cada uma contendo N bits e é considerado uma cadeia de bits
        de tamanho $L \times N$. Após os primeiros 1.000.000 bits podem vir a serem desconsiderados entào $\mathit{H_{\text{bitstring}}}$ é calculada. Tendo como amostra estimada como N x H$_{\text{bitstring}}$.
        A estimativa da entropia inicial da fonte é calculada como $\mathit{H_{\text{initial}}} = \min(\mathit{H_{\text{original}}}, \mathit{H_{\text{bitstring}}})$.

    \section{Definição de Aleatoriedade e Tipos de Geradores Aleatórios}
        RNG (Random Number Generator) é um dispositivo ou algoritmo que gera uma sequência de números ou símbolos que não possuem nenhum padrão discernível, ou seja, são imprevisíveis.
        Que utiliza uma fonte de incerteza(ou seja, a fonte de entropia) ao mesmo tempo que faz uso de uma função de processo (processo de destilação) para gerar aleatoriedade.
        O processo de destilação é uma etapa crucial para exceder qualquer ruptura de segurança na origem da entropia que possa a resultar em uma sequencia não aleatória.
        Como a geração de longas sucessoes de zeros ou uns, ou mesmo a repetição de um mesmo número, o que pode ser um indicativo de uma fonte de entropia comprometida.
        Normalmente a fonte de entropia é atraves de um processo físico, como o ruído térmico, circuito eletrico, o movimento do mouse ou a digitação do teclado, ou os efeitos quanticos de um semicondutor%\cite{NIST 800-22, pg 14. 1.1.3 RNGs}.
        Multiplas fontes podem ser utilizadas para esse processo. As saidas de um RNG podem vir a ser utilizadas para um pseudo gerador de números aleatórios (PRNG). 
        A saida de um RNG deve seguir certos critérios de aleatoriedade medidos por testes estatísticos, como os definidos pelo NIST 800--22, 
        para garantir que os números gerados sejam adequadamente aleatórios para uso em criptografia e outras aplicações sensíveis à segurança.
        Para fins de criptografia, é crucial que os números gerados por um RNG sejam verdadeiramente aleatórios. No entanto, algumas das fontes como data/hora são previsiveis, 
        que combinados com outras fontes de entropia, podem mitigar a previsibilidade. E que ainda pode se tornar deficiente ao passar pelos testes estastisticos. A produção de RNGs de alta qualidade se tornam um recurso oneroso.
        Tornando improdutiva quando se faz necessario de multiplas gerações de RNGs. 
        Por isso, a utilização de PRNGs se torna uma alternativa mais viável, pois eles podem gerar grandes quantidades de números pseudoaleatórios a partir de uma semente inicial.

        PRNG (PSeudo-Random Number Generator) é um tipo de gerador de números aleatórios que utiliza um algoritmo determinístico para produzir uma sequência de números que aparenta ser aleatória.
        Ele faz uso de uma semente (seed) inicial para gerar uma sequência de números pseudoaleatórios. Em momentos que se faz necessário uso de impreviabilidade, a seed se faz necessaria ser imprevisivel e oriundas e suma de um RNG.\@
        Por padrão um PRNG requer majoritariamente de um RNG.\@ Como o PRNG é um algoritmo determinístico, por isso ele leva o nome de pseudoaleatório, pois a sequência de números gerada por um PRNG é completamente determinada pela semente inicial.
        Como a reprodução de um sequencia de números pseudoaleatórios é possível se a semente for conhecida, somente ela se faz necessaria salvar e proteger. 
        De forma cômica, muitos dos valores gerados por um PRNG aparentam sem mais aleatórios do que os números gerados por um RNG.\@ Se sua concepção for bem desenvolvida, cada geração utiliza o valor da anterior por meio de transformações matemáticas, 
        que acabam introduzindo uma aleatoriedade adicional. Uma série de de transformações podem eliminar autocorrelações e padrões, entre a entrada e a saida.%\cite{NIST 800-22, pg 14. 1.1.3 RNGs}

        DRBG (Deterministic Random Bit Generator) é um tipo específico de PRNG que é projetado para ser criptograficamente seguro. Ele produz uma sequencia de bits que a partir de uma seed inicial, determina uma sequencia de bits pseudoaleatórios. 
        A seed para instanciar um DRBG deve ser gerada a partir de uma fonte de entropia confiável, como um RBG (Random Bit Generator), para garantir a segurança do DRBG.\@ Se a seed for protegida e tiver alta entropia, o DRBG pode ser considerado seguro para uso em criptografia.
        Além da dependência de uma semente de alta entropia, a segurança de um DRBG reside em sua capacidade de manter a Backtracking Resistance e a Prediction Resistance. 
        Conforme o NIST SP 800-90A, esses mecanismos garantem que o comprometimento temporário do estado interno do gerador não resulte na exposição de chaves criptográficas geradas anteriormente ou no colapso da imprevisibilidade dos bits futuros

        NRBG (Non-deterministic Random Bit Generator) é um tipo de gerador de números aleatórios que se baseia em processos físicos ou eventos imprevisíveis para gerar números aleatórios.
        Ao contrário dos DRBGs, que são determinísticos e dependem de uma semente para gerar números pseudoaleatórios, os NRBGs não são determinísticos e não dependem de uma semente. 
        Eles podem usar fontes de entropia, como ruído térmico, fenômenos quânticos ou outros eventos físicos imprevisíveis para gerar números aleatórios.

        TRNG (True Random Number Generator) utiliza da entropia de eventos fisicos imprevisíveis para gerar números aleatórios ao inves da utilização de algoritmos determinísticos. 
        É normalmente utilizado para gerar chaves criptográficas, Algoritmo de assinatura digital de curva eliptica[ECDSA] ou chaves MAC simetricas, para armazenar e transmitir dados com alta segurança em protocolos criptograficos como o TLS.\@
        As chaves geradas por um TRNG são altamente resistentes a ataques de força bruta, pois a entropia de um TRNG é geralmente muito alta, tornando difícil para um atacante adivinhar ou reproduzir as chaves geradas.

    \section{Avaliação Estatística da Aleatoriedade}
        
        Existem diversos testes estastisticos que podem ser utilzados em uma sequencia para avaliar e comparar com uma sequencia verdadeiramente aleatoria. 
        A aleatoriedade é uma caracteristica probabilistica, de forma que é possivel descrever em termos de probabilidade. O resultado dos avaliadores estatisticos, quando baseados em uma sequencia de bits verdadeiramente aleatoria, é denominada de \boldmath{priori} %\nocite{NIST SP 800-22, pg 16. 1.1.5 Testing}
        É capaz de aplicar uma inifinidade de testes estatísticos para avaliar a aleatoriedade de uma sequência de bits. Cada um avaliando um aspecto diferente da aleatoriedade, 
        sendo a presença ou ausencia de padrões, que sendo identificados seria uma indicação que a sequencia não possui as caracteristicas aleatórias suficientes.
        Já que existe uma possibilidade de aplicar uma sequencia infinita de testes, nenhum pode ser considerado completo, sendo necessário estabelecer um conjunto de testes estatísticos padronizados para avaliar a aleatoriedade de sequências de bits. E ainda assim devem ser analisados
        com cuidado para evitar conclusões errôneas sobre a aleatoriedade de uma sequência de bits.

        Uma hipótese nula (H$_0$) é formulada para cada teste estatístico, que geralmente afirma que a sequência de bits é aleatória. 
        Na outra ponta está a hipótese alternativa (H$_1$), que afirma que a sequência de bits não é aleatória. 
        O teste estatístico é então aplicado à sequência de bits, e um valor de teste é calculado com base nos resultados do teste. Ao final de cada teste se tem uma decisão a ser tomada, que é se rejeitar ou não a hipótese nula.

        Para cada distribuição teorica em cima da hipotese nula (H$_0$) é determinada por meios matematicos. A partir disso, na distribuição de referencia, 
        um valor critico é determinado longe o suficiente para que a probabilidade de obter um resultado tão extremo 0.9 quanto o observado seja menor do que um nivel de significancia predefinido 0.01.

        Na realização, se um valor de um teste exceder o valor crítico, a hipótese nula (H$_0$) é rejeitada. Se o pressuposto da aleatoriedade for verdadeira o valor final da estatistica téra um probabilidade muito baixa de exceder o valor critico.

        Só há duas alternativas de decisão, ou a hipótese nula (H$_0$) é rejeitada ou não é rejeitada.\@ Se a hipótese nula for rejeitada, isso sugere que a sequência de bits não é aleatória, e que pode haver padrões ou correlações presentes na sequência. 
        Se a hipótese nula (H$_0$) não for rejeitada, isso sugere que a sequência de bits é aleatória, mas não prova que a sequência é verdadeiramente aleatória.

        \boldmath{Erro Tipo 1} ocorre quando a hipótese nula (H$_0$) é rejeitada, mesmo que seja verdadeira. Isso ocorre quando o teste estatístico indica que a sequência de bits não é aleatória. A probabilidade de um erro do tipo 1 é chamaada de \boldmath{nivel de significancia} e é denotada por $\alpha$ %\nocite{NIST SP 800-22, pg 16. 1.1.5 Testing}.
        Na criptografia valores comuns para $\alpha$ são 0.01 ou 0.001, o que significa que há uma probabilidade de 1\% ou 0.1\% de rejeitar a hipótese nula (H$_0$) quando ela é verdadeira. É esperado que 1 a cada 100 sequencias de bits sejam rejeitadas.

        \boldmath{Erro do Tipo 2} ocorre quando a hipótese nula (H$_0$) não é rejeitada, mesmo que seja falsa. Isso ocorre quando o teste estatístico indica que a sequência de bits é aleatória, mesmo que haja padrões ou correlações presentes na sequência. A probabilidade de um erro do tipo 2 é denotada por $\beta$ %\nocite{NIST SP 800-22, pg 16. 1.1.5 Testing}.
        Não há um valor fixo para $\beta$, pois depende do teste estatístico específico e da natureza dos padrões ou correlações presentes na sequência de bits e que possui uma infinidade ded maneiras de que uma sequencia de bits não seja aleatória.
    
        Uma das principais metas da execução de testes estatísticos é minimizar a probabilidade de cometer erros do tipo 2. Diminuindo a probabilidade de ter um gerador ruim aceitar uma sequência de bits como aleatória.
        O estatístico de teste é usado para calcular um P-value que resume a força da evidência contra a hipótese nula.
        
        Se P-value = 1, a sequência parece ter aleatoriedade perfeita.
        
        Se P-value = 0, a sequência parece ser completamente não-aleatória.
        
        Se P-value $\geq \alpha$, a hipótese nula é aceita (a sequência parece ser aleatória).
        
        Se P-value $< \alpha$, a hipótese nula é rejeitada (a sequência não é aleatória).

        O NIST SP 800--22 é um conjunto de testes estatísticos amplamente utilizado para avaliar a aleatoriedade de sequências de bits geradas por geradores de números aleatórios. 
        Sendo ao todo 15 testes, cada um projetado para detectar diferentes tipos de padrões ou falhas na aleatoriedade. Esses testes incluem:
         \begin{enumerate}
          \item Teste de Frequência (monobit)
          \item Teste de Frequência dentro de um bloco
          \item Teste de Sequencias
          \item Testes para Maior Sequencia de Uns em um bloco
          \item Testes de Posto de Matriz Binária
          \item Teste da Transformada Discreta de Fourier (Espectral)
          \item Teste de Correspondência de Modelos não Sobrepostos
          \item Teste de Correspondência de Modelos Sobrepostos
          \item Teste "Estatístico Universal" de Maurer
          \item Teste de Complexidade Linear
          \item Teste Serial
          \item Teste de Entropia Aproximada
          \item Teste de Somas Cumulativas (Cusums)
          \item Teste de Excursões Aleatórias
          \item Teste de Variante de Excursões Aleatórias
        \end{enumerate}%\nocite{NistsTP}

        %\nocite{NIST SP 800-22, pg 16. 1.1.5 Testing}
        